\begin{definition}[Native dual-path words]
    \label{def:peps_dual_path_word}
    \lean{TNLean.PEPS.torusDualStepWord,TNLean.PEPS.torusDualPathWord}\leanok
    Let $X$ be the finite native vertex set of a torus with positive periods.
    The free group on $X\sqcup X$ has one east and one north generator at
    each actual native base vertex. West and south steps use their inverses.
    Write $W(p)$ for the reverse-ordered word of a labelled path, so
    $W(pq)=W(q)W(p)$ and $W(p^{-1})=W(p)^{-1}$.
\end{definition}

\begin{definition}[Fixed native root routes]
    \label{def:peps_dual_root_routes}
    \lean{TNLean.PEPS.torusDualRootRoute,TNLean.PEPS.torusDualRootRoute_zero}\leanok
    Fix the origin $o=(0,0)$. Assign to each vertex $v$ the actual native
    comb route $r_v:o\to v$, using its least nonnegative integer coordinates.
    Assign $r_o$ to be the constant path.
\end{definition}

\begin{definition}[Finite presentation for allowed dual-square supports]
    \label{def:peps_dual_finite_presentation}
    \lean{TNLean.PEPS.torusDualFaceWord,TNLean.PEPS.torusDualPresentationRelators,
        TNLean.PEPS.card_torusDualPresentationRelators_le}\leanok
    For allowed primal supports $U\subseteq X$, use the relator $W(r_v)$
    at each vertex and the relator $W(\mathrm{EN})W(\mathrm{NE})^{-1}$ for
    the square enclosing each $c\in U$. Let $P_U$ be the quotient of the
    free group by the normal closure of these relators, and let $\pi_U$
    be its quotient map. There are two generators per native vertex and
    at most $2|X|$ defining relators. Labels remain distinct at periods one and two.
\end{definition}

\begin{lemma}[Classification by the finite native presentation]
    \label{lem:peps_dual_finite_presentation_classification}
    \lean{TNLean.PEPS.TorusDualHomotopy.presentation_word_eq,
        TNLean.PEPS.torusDualHomotopy_iff_presentation_word_eq}\leanok
    For any two finite native paths $p,q$ with the same endpoints,
    generated homotopy supported on $U$ is equivalent to
    $\pi_U(W(p))=\pi_U(W(q))$ in $P_U$.
    This is equality in one explicit finitely presented group; no terminating
    word-normalization algorithm or finite-group test family is asserted.
\end{lemma}
\begin{proof}\leanok
    \uses{def:peps_dual_path_word,def:peps_dual_root_routes,
        def:peps_dual_finite_presentation,lem:peps_dual_flat_transport_completeness}
    Backtracks cancel and the allowed square words are defining relators,
    so every generated homotopy preserves the quotient word.
    Conversely, take any reversible fibre transport flat on the allowed
    squares. Its root routes give bijections $T_v:Y_o\to Y_v$.
    Normalize each step $E_e:Y_a\to Y_b$ to the root-fibre permutation
    $T_b^{-1}E_eT_a$. An actual path then has normalized action
    $T_b^{-1}E_pT_a$. Each fixed root-route word acts as identity, and
    local flatness makes every allowed square relator act as identity.
    Thus these permutations define a homomorphism from $P_U$.
    Equal quotient words give equal normalized path actions; undoing the
    root-route bijections gives equal actions of every locally flat transport.
    The flat-transport classification yields the required supported homotopy.
\end{proof}
